\clearpage
\phantomsection% 
\addcontentsline{toc}{chapter}{ABSTRACT}
\thispagestyle{fancy}

\begin{center}
	\large \bfseries \MakeUppercase{Abstract}\\
	\normalsize \normalfont {\thetitleEN}\\
	\normalsize \normalfont {\theauthor}\\
	\bigskip
	
	\normalsize \normalfont \justifying \singlespacing
	Falls among older adults are a serious problem that requires special attention to prevent injuries. According to the World Health Organization (WHO), about 37.3 million fall incidents require medical attention, and around 684,000 deaths worldwide are caused by falls. In addition, older adults living in nursing homes face a higher risk of falling than those living at home, with an annual incidence of 30--50\% and about 40\% experiencing recurrent falls. This issue calls for a fast and accurate fall-position detection model that can operate continuously for 24 hours, especially in areas that are not always supervised, to improve older adults' quality of life. Previous studies developed vision-based fall detection models using CNN and LSTM, but a key limitation is the high computational power required for image processing. This study aims to develop a detection model based on Graph Convolutional Networks (GCN) using pose landmark data, which is technically lighter than image data. In this work, a GCN model was developed with two hyperparameter configurations, namely the default configuration and a configuration derived from an ablation study. Experimental results show that the ablation-study configuration achieves a significant performance improvement over the default configuration. Overall (average), accuracy increases by 0.018, precision by 0.024, F1-score by 0.049, and the largest gain appears in recall, by 0.093. These results indicate that the ablation-study configuration can improve the performance of the GCN model in detecting falls among older adults, with strong potential for application in future elderly health monitoring systems focused on fall and non-fall events.
    \vspace{20pt}
	
	\textbf{Keywords: Eldery, Fall, Deep Learning, Graph Convolution Network, Mediapipe, Pose Landmark, Ablation Study}
	
	\vfill
	
\end{center}
\clearpage